\documentclass[11pt]{article}
\usepackage[margin=1in]{geometry}
\usepackage[utf8]{inputenc}
\usepackage[T1]{fontenc}
\usepackage{lmodern}
\usepackage{amsmath,amssymb,amsthm}
\usepackage{booktabs}
\usepackage[title]{appendix}
\usepackage[colorlinks=true,linkcolor=blue,citecolor=blue,urlcolor=blue,pdfauthor={P. Graham},pdftitle={Cyclic sieving for q-Kreweras numbers in the exceptional types}]{hyperref}
\newtheorem{theorem}{Theorem}[section]
\newtheorem{corollary}[theorem]{Corollary}
\newtheorem{proposition}[theorem]{Proposition}
\newtheorem{lemma}[theorem]{Lemma}
\theoremstyle{definition}
\newtheorem{example}[theorem]{Example}
\newcommand{\NC}{\mathrm{NC}}
\newcommand{\Krew}{\mathrm{Krew}}
\newcommand{\Fix}{\mathrm{Fix}}
\newcommand{\Z}{\mathbb{Z}}
\newcommand{\Oc}{\mathcal{O}}
\setlength{\emergencystretch}{3em}
\title{Cyclic sieving for $q$-Kreweras numbers in the exceptional types}
\author{P. Graham\\ \small\url{https://orcid.org/0009-0007-8493-1505}}
\date{November 4, 2026}
\begin{document}
\maketitle
\begin{abstract}
Let $W$ be a finite Weyl group with Coxeter number $h$ and simple Lie
algebra $\mathfrak g$. For each $s\ge1$ the cyclic group $\Z/sh\Z$ acts on
the multichains $w_1\le\dots\le w_s$ of noncrossing partitions of $W$.
Reiner and Sommers defined, for each nilpotent orbit $\Oc$ of $\mathfrak g$
and each positive integer $m$, a $q$-Kreweras number, a polynomial in $q$,
and tabulated it in the exceptional types. Each $W$-orbit of flats $X$
determines a nilpotent orbit $\Oc_X$, principal in the Levi subalgebra of
$X$. Their Conjecture~1.4 says that the $q$-Kreweras number of $\Oc_X$ with
$m=sh+1$, at a primitive $d$-th root of unity, counts the multichains fixed
by an element of order $d$ whose bottom element has fixed space in the
orbit of $X$. Reiner and Sommers proved it in types $A$, $B$, $C$ and $D$.
We prove it in types $G_2$, $F_4$, $E_6$, $E_7$ and $E_8$, so it holds for
every irreducible Weyl group. At $q=1$ the statement is a count proved by
Douvropoulos and Josuat-Vergès; what is new is every $d>1$. Krattenthaler
and Müller proved the unrefined cyclic sieving, Armstrong's conjecture, in
the exceptional types by reducing it to finitely many cases. Using their
parametrization of $s$ and of the acting element by three integers, and
their description of the fixed points, we show that with two of the
integers fixed the fixed-point count is a polynomial of degree at most the
rank in the third. A finite-difference argument shows the same for the
tabulated $q$-Kreweras number at the root of unity, once it holds for rank
plus one consecutive values of the third integer. Only finitely many choices of the fixed pair need a
direct check. The resulting finitely many identities are verified by two
independent exact programs, and the reduction is checked in Lean.

This work has not yet been reviewed by experts.
The proofs, the Lean formalization, the computer checks and the text were
produced by AI models (from OpenAI and Anthropic), working under the
author's direction.
\end{abstract}
\medskip
\noindent\textit{Keywords:} cyclic sieving, noncrossing partitions,
Kreweras numbers, Weyl groups, nilpotent orbits, Fuss--Catalan
combinatorics.

\noindent\textit{MSC 2020:} 05E10, 05A15, 20F55, 17B08.

\section{Introduction}

\subsection{Main result}\label{sec:intro-main}
Let $\Phi$ be an irreducible crystallographic root system of rank $n$ in a
real vector space $V$, with Weyl group $W$ and Coxeter number $h$. The \emph{absolute length} $\ell(w)$ of $w\in W$ is the
least number of reflections whose product is $w$; it equals
$n-\dim V^w$, where $V^w$ is the space of vectors fixed by $w$
\cite[Lemma~2]{Car}. The \emph{absolute order} is $u\le w$ if
$\ell(u)+\ell(u^{-1}w)=\ell(w)$. Fix a Coxeter element $c$, a product of
the simple reflections in some order. The noncrossing partitions are
$\NC(W)=\{w\in W: w\le c\}$ \cite{Bes} (all Coxeter elements are conjugate,
and conjugation is an automorphism of the absolute order, so the choice of
$c$ does not matter), and $\NC^{(s)}(W)$ is the set of multichains
$w_1\le w_2\le\dots\le w_s$ in $\NC(W)$ \cite{Arm}. Such a multichain is the same as a \emph{delta
sequence} $(\delta_0;\delta_1,\dots,\delta_s)$ with
$\delta_0\delta_1\cdots\delta_s=c$ and
$\ell(\delta_0)+\dots+\ell(\delta_s)=n$, through $\delta_0=w_1$ and
$\delta_i=w_i^{-1}w_{i+1}$ ($w_{s+1}=c$). The cyclic group $\Z/sh\Z$
acts on delta sequences through the generator of Armstrong's action, in
the form of Bessis and Reiner \cite[(6.4)]{BR}:
\begin{equation}\label{eq:phi}
\varphi(\delta_0;\delta_1,\dots,\delta_s)
=(v;\,c\delta_sc^{-1},\delta_1,\dots,\delta_{s-1}),\qquad
v=(c\delta_sc^{-1})\,\delta_0\,(c\delta_sc^{-1})^{-1}.
\end{equation}
The new bottom element $v$ is conjugate to $\delta_0$, so the $W$-orbit of
the fixed space of the bottom element $w_1=\delta_0$ is constant on
$\Z/sh\Z$-orbits.

A \emph{flat} is an intersection of reflecting hyperplanes; each $V^w$ is
one, and $[X]=W\cdot X$ denotes the $W$-orbit of a flat $X$. Let
$\mathfrak g$ be the complex simple Lie algebra of $\Phi$ with Cartan
subalgebra $\mathfrak h$. The centralizer of $X\subset\mathfrak h$ is a
Levi subalgebra $\mathfrak l_X$, and $\Oc_X$ is the nilpotent orbit
containing the principal nilpotent elements of $\mathfrak l_X$
\cite[\S1.3]{RS}. For a nilpotent orbit $\Oc$ and a positive integer $m$,
Reiner and Sommers define the $q$-Kreweras number
$\Krew(\Phi,\Oc,m;q)=f_{e,1}(m;q)$, a polynomial in $q$ that they attach to
an element $e\in\Oc$ \cite[\S1.3]{RS}; the conjecture uses $m=sh+1$. They
list these polynomials in the exceptional types in the tables of
\cite[\S3.6]{RS}, and throughout $\Krew(\Phi,\Oc_X,m;q)$ denotes the formula
tabulated there. Write $[x]_q=(1-q^x)/(1-q)$.

\begin{theorem}\label{thm:main}
Let $W$ be of type $G_2$, $F_4$, $E_6$, $E_7$ or $E_8$. Let $s\ge1$, let
$d$ divide $sh$, let $\omega$ be a primitive $d$-th root of unity, and let
$[X]$ be a $W$-orbit of flats. Then $\Krew(\Phi,\Oc_X,sh+1;q=\omega)$ is the
number of multichains $w_1\le\dots\le w_s$ in $\NC^{(s)}(W)$ with
$V^{w_1}\in[X]$ that are fixed by an element of order $d$ of $\Z/sh\Z$.
\end{theorem}

In the language of Reiner, Stanton and White \cite{RSW}, the theorem says
that for each $[X]$ the multichains in $\NC^{(s)}(W)$ with
$V^{w_1}\in[X]$, the polynomial $\Krew(\Phi,\Oc_X,sh+1;q)$ and the cyclic
group $\Z/sh\Z$ exhibit the cyclic sieving phenomenon: the polynomial's
value at a primitive $d$-th root of unity counts the points fixed by an
element of order $d$.

\begin{corollary}\label{cor:all}
Conjecture~1.4 of Reiner and Sommers \cite{RS} holds for every irreducible
Weyl group.
\end{corollary}

\begin{proof}
Reiner and Sommers prove it in types $A$, $B$, $C$, $D$ \cite[Theorem~1.7]{RS}.
In the other types they compute $f_{e,1}$ in the tables of \cite[\S3.6]{RS},
which give the function $\Krew$ of Theorem~\ref{thm:main}, so
Theorem~\ref{thm:main} proves the conjecture there.
\end{proof}

\begin{example}\label{ex:g2}
Take $W$ of type $G_2$ ($n=2$, $h=6$) and $s=1$, so $m=7$. Then the
generator of $\Z/6\Z$ acts on $\NC(W)$ by conjugation by $c$, an element of
order $6$; since $c^3=-1$ is central, $c^3$ acts trivially. $\NC(W)$
consists of the identity (the flat $V$), three reflections of each root
length and $c$ (the flat $\{0\}$). The four orbits of flats are labeled by
the Bala--Carter labels $0$, $A_1$, $\tilde A_1$ and $G_2$ of $\Oc_X$, and
the tables of \cite[\S3.6]{RS} give
\[
\tfrac{[m-1]_q[m-5]_q}{[2]_q[6]_q},\quad
q^{m-5}\tfrac{[m-1]_q}{[2]_q},\quad
q^{m-3}\tfrac{[m-1]_q}{[2]_q},\quad
q^{2m-2},
\]
which at $m=7$ are $1$, $q^2(1+q^2+q^4)$, $q^4(1+q^2+q^4)$ and $q^{12}$.
Their values at a primitive $d$-th root of unity are:
\begin{center}
\begin{tabular}{lrrrr}
\toprule
Label $\backslash$ $d$ (element) & $1$ ($1$) & $2$ ($c^3$) & $3$ ($c^2$) & $6$ ($c$)\\
\midrule
$0$ (the identity) & 1 & 1 & 1 & 1\\
$A_1$ (long reflections) & 3 & 3 & 0 & 0\\
$\tilde A_1$ (short reflections) & 3 & 3 & 0 & 0\\
$G_2$ ($c$) & 1 & 1 & 1 & 1\\
\bottomrule
\end{tabular}
\end{center}
Each entry is both the polynomial's value and the number of elements of
$\NC(W)$ with that label fixed by conjugation by the element in the column.
\end{example}

The Kreweras numbers refine the Fuss--Catalan number $|\NC^{(s)}(W)|$
according to the $W$-orbit of the fixed space of the bottom element
\cite[\S1.1]{RS}. Reiner and Sommers' $q$-analogues are attached to
nilpotent orbits of $\mathfrak g$ \cite[\S1.3]{RS}, so the conjecture links
the nilpotent orbits of the Lie algebra with the cyclic action on the
noncrossing partitions of its Weyl group. With Theorem~\ref{thm:main} it
holds uniformly for every Weyl group.

\subsection{Prior work}

Reiner and Sommers stated the conjecture in 2016 \cite{RS}, answering in
conjectural form a question that Bessis and Reiner posed for $s=1$
\cite[\S6.1, Question~6.3]{BR}. Their Conjecture~1.4 reads: ``For $m=sh+1$, and for each $W$-orbit $[X]$,
$\Krew(\Phi,\Oc_X,m;q=\omega_d)$ counts those $w_1\le\dots\le w_s\in
\NC^{(s)}(W)$ which are fixed under the action an element of order $d$ in
the $\Z/sh\Z$-action and have $V^{w_1}\in[X]$.'' Their Theorem~1.7 proves it
in types $A$, $B$, $C$ and $D$. They describe the $q=1$ count as follows:
``In classical ($A,B,C,D$) and dihedral ($I$) types, work of Rhoades showed
more generally that $\Krew(W,[X],sh+1)$ counts the elements
$w_1\le\dots\le w_s$ in $\NC^{(s)}(W)$ with $V^{w_1}$ in $[X]$; this remains
open in the exceptional types $E,F,H$'' \cite[\S1.1]{RS}. Douvropoulos and
Josuat-Vergès proved this count for every irreducible $W$
\cite[Theorem~3.10]{DJV}, a theorem they attribute to Rhoades \cite{Rho}
except for $E_7$ and $E_8$; for the remaining exceptional groups they write
that ``this can be done by computer''. When $e\in\Oc_X$,
$f_{e,1}(m;1)$ is this Kreweras number \cite[\S1.3]{RS}, so the
case $d=1$ of Theorem~\ref{thm:main} is their theorem, and what is new here
is every $d>1$. The case $s=1$ of the count is a theorem of Athanasiadis
and Reiner \cite[Theorem~6.3]{AR}, as Reiner and Sommers recall
\cite[\S1.1]{RS}.

Without the condition on $V^{w_1}$, the analogous statement for the
$q$-Fuss--Catalan number (a product of ratios of $q$-integers over the
degrees of $W$) is Armstrong's Conjecture~5.4.7 \cite{Arm}. Reiner and
Sommers record it as their Theorem~1.3, citing Bessis and Reiner
\cite{BR} \cite[\S1.2]{RS}. Krattenthaler
and Müller prove that the generalized noncrossing partitions of the
well-generated complex reflection groups of exceptional type ``obey two different cyclic sieving
phenomena, as conjectured by Armstrong, and by Bessis and Reiner''
\cite{KM}. To do so they ``reduce the conjectures for
each of the 26 exceptional types to a finite problem'' \cite[\S1]{KM}. Our
proof follows that reduction: it uses their formula \cite[(5.1)]{KM} for
powers of $\varphi$, the parametrization and the block equations
\cite[(5.8)]{KM} from the proof of their Lemma~30, and their Lemma~31 for
large $\nu$.

Three works treat related refinements that Theorem~\ref{thm:main} does
not address. Krattenthaler and Stump pose as an open problem a cyclic
sieving phenomenon for positive $m$-divisible noncrossing partitions with
prescribed block structure in exceptional types
\cite[\S7.2, Question~7.5 and Open Problem~7.6]{KS}. Pouillart provides ``a
uniform refinement of these cyclic sieving phenomena using a $q$-analogue''
of the face formula of Douvropoulos and Josuat-Vergès, for faces of the
generalized cluster complex under the Fomin--Reading rotation \cite{Pou}.
Kim defines, ``for a coincidental Coxeter group, i.e.\ of type $A_{n-1},
BC_n, H_3$, or $I_2(m)$'', $q$-Kreweras numbers attached to limit symbols
\cite{Kim}. Our result concerns neither the positive objects, nor faces, nor
these other $q$-analogues.

\subsection{Method}
The conjecture is a statement for infinitely many $s$. Following
Krattenthaler and Müller, we write $s=m_1\nu$ and the power of $\varphi$
as $\varphi^{m_1h_1}$ with $h_1\mid h$ and $\gcd(h_1,\nu)=1$. For fixed
$(h_1,\nu)$ the fixed points with a given $[X]$ are counted by a polynomial
in $m_1$ of degree at most $n$, built from chains in the centralizer of
$c^{h_1}$, as in the proof of \cite[Lemma~30]{KM}
(Proposition~\ref{prop:fix}). The new step for the refinement is the
$q$-Kreweras side: the tabulated $q$-Kreweras numbers are quotients of
products and sums of $q$-integers, with no product formula of the kind used
for the $q$-Fuss--Catalan numbers. We show directly from the shape of each
formula that if its value at $\omega$ is regular and agrees with a
polynomial in $m_1$ of degree at most $n$ for $n+1$ consecutive values of
$m_1$, then it does so for every $m_1$ (Proposition~\ref{prop:rhs}). For
$\nu>n$ both sides are given in closed form outside finitely many pairs
$(h_1,\nu)$. The remaining identities, $104$ pairs $(h_1,\nu)$ with $\nu\le
n$ and $21$ with $\nu>n$ over the five types, are checked by computer
(Section~\ref{sec:comp}), and the reduction is formalized in Lean
(Section~\ref{sec:lean}).

\section{The reduction to finitely many identities}\label{sec:red}

Write $\Fix_X(s,p)$ for the number of delta sequences of length $s$ fixed by
$\varphi^p$ whose bottom element $\delta_0$ has $V^{\delta_0}\in[X]$, and
$R_X(m;q)=\Krew(\Phi,\Oc_X,m;q)$.

\paragraph{Step 1: the parameters.}
An element of order $d$ in $\Z/sh\Z$ generates the same subgroup as
$\varphi^p$ with $p=sh/d$, so it has the same fixed points.

\begin{lemma}\label{lem:param}
The pairs $(s,p)$ with $s\ge1$ and $p\mid sh$ correspond one to one to
the triples $(h_1,\nu,m_1)$ of positive integers with $h_1\mid h$ and
$\gcd(h_1,\nu)=1$, through $s=m_1\nu$ and $p=m_1h_1$. Then $d=sh/p=\nu h/h_1$.
\end{lemma}

\begin{proof}
Given $(s,p)$, put $m_1=\gcd(p,s)$, $\nu=s/m_1$ and $h_1=p/m_1$. Then
$\gcd(h_1,\nu)=1$, and $p\mid sh$ gives $h_1\mid \nu h$, hence $h_1\mid h$.
Conversely $s=m_1\nu$ and $p=m_1h_1$ satisfy $p\mid sh$ and
$\gcd(p,s)=m_1$.
\end{proof}

\paragraph{Step 2: the fixed points.}
Fix $h_1\mid h$ and $\nu\ge1$ with $\gcd(h_1,\nu)=1$, and write
$h_1=t\nu+r$ with $0\le r<\nu$, as in the proof of \cite[Lemma~30]{KM}.
Since $\gcd(r,\nu)=1$, the map sending $k$ to $k+r$, or to $k+r-\nu$ when
$k+r>\nu$, is one cycle on $\{1,\dots,\nu\}$. Starting from $e_1=0$, follow
this cycle for $\nu-1$ steps, setting $e_{k+r}=e_k+t$ for a step with
$k+r\le\nu$ and $e_{k+r-\nu}=e_k+t+1$ for a step with $k+r>\nu$. This
defines $e_1,\dots,e_\nu$. Going once round the cycle would add
$\nu t+r=h_1$, so the step that closes it defines nothing; in the proof
below it becomes the condition of commuting with $c^{h_1}$. Equivalently,
$e_k$ is the unique integer in $[0,h_1)$ with $h_1\mid \nu e_k+(k-1)$, and
for elements commuting with $c^{h_1}$ any exponents with this divisibility
give the same conjugates below (this is the form used in
Appendix~\ref{app:lean}). Let $C$ be the centralizer of $c^{h_1}$ in
$W$. For $u\in\NC(W)\cap C$ put $u_k=c^{e_k}uc^{-e_k}$ and
$\delta_0(u)=c\,(u_1u_2\cdots u_{\nu})^{-1}$, and let
\[
U_X=\{u\in\NC(W)\cap C:\ \ell(\delta_0(u))+\nu\,\ell(u)=n,\ V^{\delta_0(u)}\in[X]\}.
\]
For $u\in\NC(W)$ let $N_k(u)$ be the number of chains
$1=z_0<z_1<\dots<z_k=u$ in the absolute order with every $z_i\in C$.

\begin{proposition}\label{prop:fix}
For every $m_1\ge1$,
\[
\Fix_X(m_1\nu,\,m_1h_1)=L_X(m_1):=\sum_{u\in U_X}\ \sum_{k=0}^{n}\binom{m_1}{k}N_k(u).
\]
In particular $L_X$ is a polynomial in $m_1$ of degree at most $n$.
\end{proposition}

\begin{proof}
Put $s=m_1\nu$ and $p=m_1h_1=ts+m_1r$. On the last $s$ entries,
$\varphi$ acts by $\rho(\delta_1,\dots,\delta_s)=(c\delta_sc^{-1},\delta_1,
\dots,\delta_{s-1})$, and $\delta_0=c(\delta_1\cdots\delta_s)^{-1}$ is
determined by them, so the fixed points of $\varphi^p$ are those of
$\rho^p$. Since $\rho^s$ is conjugation by $c$ in each entry,
$\rho^p(\delta)_i=c^{t+1}\delta_{i-m_1r+s}c^{-t-1}$ for $i\le m_1r$ and
$\rho^p(\delta)_i=c^t\delta_{i-m_1r}c^{-t}$ for $i>m_1r$
\cite[(5.1), (5.8)]{KM}. Cut $(\delta_1,\dots,\delta_s)$ into $\nu$ blocks
$B_1,\dots,B_{\nu}$ of $m_1$ consecutive entries. As the shift $m_1r$ is a
multiple of $m_1$, the equation $\rho^p\delta=\delta$ says that
$B_{k+r}=c^tB_kc^{-t}$ when $k+r\le \nu$ and $B_{k+r-\nu}=c^{t+1}B_kc^{-t-1}$
otherwise, entry by entry. Following the cycle from $B_1$, as in the proof
of \cite[Lemma~30]{KM}, these equations hold exactly when every entry of $B_1$ commutes with $c^{h_1}$ and
$B_k=c^{e_k}B_1c^{-e_k}$ for all $k$.

Let $u$ be the product of the entries of $B_1$. The product of the entries
of $B_k$ is then $u_k$, so $\delta_0=\delta_0(u)$, and by subadditivity and
conjugation invariance of $\ell$,
\[
n=\ell(\delta_0)+\nu\sum_{x\in B_1}\ell(x)\ \ge\ \ell(\delta_0)+\nu\,\ell(u)
=\ell(\delta_0)+\sum_k\ell(u_k)\ \ge\ \ell(c)=n.
\]
Both inequalities are equalities: $B_1$ is a length-additive factorization
of $u$ into $m_1$ elements of $C$, and
$\ell(\delta_0(u))+\nu\ell(u)=n$. The second condition makes
$c=\delta_0u_1\cdots u_{\nu}$ length-additive, so $u=u_1\le c$ and $u\in
\NC(W)\cap C$. Conversely, each $u\in\NC(W)\cap C$ with
$\ell(\delta_0(u))+\nu\ell(u)=n$ and each length-additive factorization of
$u$ into $m_1$ elements of $C$ gives a fixed point. Length-additive
factorizations $u=x_1\cdots x_{m_1}$ correspond to multichains
$1=y_0\le y_1\le\dots\le y_{m_1}=u$ through $y_j=x_1\cdots x_j$, and all
$x_j$ lie in $C$ exactly when all $y_j$ do. A multichain with exactly $k$
strict steps is a strict chain of length $k$ together with a choice of the
$k$ strict positions among $m_1$, which gives $\binom{m_1}{k}N_k(u)$.
\end{proof}

\begin{lemma}[{\cite[Lemma~31]{KM}}]\label{lem:large}
If $\nu>n$, then $L_X(m_1)=1$ when $X=\{0\}$ and $L_X(m_1)=0$ otherwise.
\end{lemma}

\begin{proof}
If $u\ne1$ then $\ell(\delta_0(u))+\nu\ell(u)\ge \nu>n$, so $U_X\subseteq\{1\}$.
For $u=1$, $\delta_0=c$, whose fixed space is $\{0\}$, and $N_k(1)$ is $1$
for $k=0$ and $0$ for $k>0$.
\end{proof}

\paragraph{Step 3: the right side.}
Every formula $R_X(m;q)$ in the tables of \cite[\S3.6]{RS} that we use is a
product of factors of four kinds: $q^{\alpha m+\beta}$ with integers
$\alpha,\beta$; numerator brackets $[m-a]_q$; sums
$\sum_j\pm q^{\gamma_j}[m-a_j]_q$ of such brackets; and denominator brackets
$1/[b]_q$ with $b\ge1$. In each tabulated formula (a row of the tables)
there are at most $n$ factors of the second and third kinds together,
and at most $n$ of the fourth kind; the programs check this for every row.

\begin{proposition}\label{prop:rhs}
Let $R(m;q)$ be a product of factors of the four kinds, with at most $n$
of the second and third kinds together and at most $n$ of the fourth kind.
Let $E\ge1$, and let $\omega$ be a root of unity with
$\omega^E=1$. Suppose that for $m_1=1,\dots,n+1$ the rational function
$R(Em_1+1;q)$ is regular at $q=\omega$ with value $P(m_1)$, where $P$ is a
polynomial of degree at most $n$. Then the same holds for every $m_1\ge1$.
\end{proposition}

\begin{proof}
Put $Z=q^{Em_1}$, so that $q^m=qZ$ for $m=Em_1+1$. Then
$q^{\alpha m+\beta}=q^{\alpha+\beta}Z^\alpha$ and
$[m-a]_q=(1-q^{1-a}Z)/(1-q)$. Multiplying out,
$R(Em_1+1;q)=\sum_i\psi_i(q)Z^{\alpha_i}$ with finitely many integers
$\alpha_i$ and rational functions $\psi_i$ that do not depend on $m_1$.
Each $\psi_i$ has a pole of order at most $n$ at $\omega$: if $\omega=1$
the poles come only from the factors $1/(1-q)$, one for each factor of the
second or third kind, and if $\omega\ne1$ only from the denominator brackets,
each of which has at most a simple zero at $\omega$. Let $\Delta$ be the
forward difference in $m_1$. Since $\Delta Z^\alpha=Z^\alpha(q^{E\alpha}-1)$,
\[
\Delta^{n+1}R(Em_1+1;q)=\sum_i\psi_i(q)\,Z^{\alpha_i}\,(q^{E\alpha_i}-1)^{n+1}.
\]
As $\omega^E=1$, each factor $(q^{E\alpha_i}-1)^{n+1}$ vanishes to order at
least $n+1$ at $\omega$, so every term is regular at $\omega$ with value $0$.
Now $R(E(m_1+n+1)+1;q)$ equals $\Delta^{n+1}R(Em_1+1;q)$ minus a fixed
integer combination of $R(E(m_1+j)+1;q)$ for $0\le j\le n$, and $P$ satisfies
the same relation with $\Delta^{n+1}P=0$. Induction on $m_1$ completes the
proof.
\end{proof}

\paragraph{Step 4: proof of Theorem~\ref{thm:main}.}
Let $s$, $d$ and $[X]$ be given, and let $(h_1,\nu,m_1)$ correspond to
$(s,sh/d)$ by Lemma~\ref{lem:param}. Put $E=\nu h$, so $sh+1=Em_1+1$ and
$\omega^E=1$ because $d=\nu h/h_1$ divides $E$. By
Proposition~\ref{prop:fix} the left side is $L_X(m_1)$, a polynomial of
degree at most $n$.

If $\nu\le n$, there are finitely many pairs $(h_1,\nu)$, and for each of
them and each $[X]$ the computation of Section~\ref{sec:comp} verifies that
$R_X(Em_1+1;q)$ is regular at $\omega$ with value $L_X(m_1)$ for
$m_1=1,\dots,n+2$. Proposition~\ref{prop:rhs} gives every $m_1$.

If $\nu>n$, then by Lemma~\ref{lem:large} $L_X$ is $1$ for $X=\{0\}$ and $0$
otherwise. Let $b_{\max}$ be
the largest $b$ of a denominator bracket in the type's table ($6$, $12$, $12$,
$18$ and $30$ for $G_2$, $F_4$, $E_6$, $E_7$, $E_8$). If $d>b_{\max}$, no
denominator bracket vanishes at $\omega$. Every row except that of
$X=\{0\}$ has a factor $[m-1]_q$, which vanishes at $\omega\ne1$ because
$\omega^{m-1}=\omega^{Em_1}=1$, and the row of $X=\{0\}$ is $q^{n(m-1)}$,
which is $1$ at $\omega$; the programs check both facts. If $d\le b_{\max}$,
only finitely many pairs $(h_1,\nu)$ occur, and the computation verifies
that $R_X(Em_1+1;\omega)$ is $1$ for $X=\{0\}$ and $0$ otherwise, for
$m_1=1,\dots,n+2$; again
Proposition~\ref{prop:rhs} gives every $m_1$. This covers every $s$, every
$d\mid sh$ and every $[X]$. \qed

\section{Computational verification}\label{sec:comp}

For each type the programs build $\Phi$, the reflections as integer
matrices and $c=s_1\cdots s_n$ in Bourbaki's numbering, with
$\ell(w)=\operatorname{rank}(w-1)$ \cite[Lemma~2]{Car}. They compute
$\NC(W)$ and check that its size is the Catalan number of $W$ ($8$, $105$,
$833$, $4160$, $25080$), and the absolute order on it, whose number of
comparable pairs $u\le w$ is checked against the Fuss--Catalan number
$\mathrm{Cat}^{(2)}(W)$.

The programs then match each $W$-orbit of flats with its row in the tables
of \cite[\S3.6]{RS} and check that the match is a bijection onto the $4$,
$12$, $17$, $32$ and $41$ rows of orbits principal in a Levi subalgebra
(Bala--Carter labels without parentheses), taken with the trivial
representation of the component group. Two invariants decide the match.
The type of the parabolic subgroup fixing the flat must be the Levi type of
the label. The number of hyperplanes of the arrangement restricted to the
flat must be the sum of its Orlik--Solomon exponents; these are the
representation exponents listed in the second column of all the rows of
the orbit \cite[\S2 after (2.2), and \S3.6]{RS}, and for ten orbits one of
them is listed in a row with a nontrivial representation. The second
invariant separates three pairs of orbits in $E_7$ with the same Levi type.
These $106$ rows were transcribed by a program from the arXiv source of
\cite{RS} and compared with the journal pages.

For each $h_1\mid h$ and $1\le \nu\le n$ with $\gcd(h_1,\nu)=1$, and for the
pairs with $\nu>n$ and $d\le b_{\max}$ (Appendix~\ref{app:cases}), a
certificate lists the coefficients of $L_X$ in the basis $\binom{m_1}{k}$ for
every $[X]$. Two checkers, written separately and sharing no code,
recompute everything from the root system and compare with the
certificate:
\begin{itemize}
\item Checker A builds $\NC(W)$ bottom-up with exact ranks, finds a standard
parabolic subgroup conjugate to the one fixing each flat, and evaluates
each tabulated formula as an exact polynomial in $\Z[q]$ (the division in
each formula is checked exact), reduced modulo the cyclotomic polynomial of
order $d$; the residue must be the integer $L_X(m_1)$.
\item Checker B lets $W$ act by permutations of the roots, builds
$\NC(W)$ top-down from $c$, identifies orbits of flats by the type of the
roots vanishing on $V^w$ and the number of restricted hyperplanes, derives
the exponents $e_k$ by iterating $\rho$ symbolically, and evaluates each
tabulated formula from the leading term of each factor's expansion around
$q=\omega$, in $\mathbb Q(e^{2\pi i/d})$.
\end{itemize}
Both check every identity for $m_1=1,\dots,n+2$ and the row shapes used in
Step~3. Checker B also recounts fixed points by brute force, by the literal
action \eqref{eq:phi} and by $\rho$, for the small $s$ given in
Appendix~\ref{app:cases}. The identities with $d=1$ are the counts of
\cite[Theorem~3.10]{DJV}, and they agree for all $s$. Deliberately
corrupted certificates and tables are rejected by both checkers. The README
gives the commands.

\section{Formal verification}\label{sec:lean}

Steps 1--4 are proved in Lean~4 \cite{Lean} with Mathlib \cite{Mathlib}.
The theorem \texttt{RS16.Main.conjecture\_1\_4} has the following form.

\emph{Data:} a finite group $G$; a function $\ell:G\to\mathbb N$; an element
$c\in G$ and integers $h$ and $n$; a map $\mathrm{cls}$ from $G$ to a set of
labels; for each label $X$ a formula $\mathrm{rows}(X)$ of the shape of
Step~3; and block exponents $e$.

\emph{Hypotheses} (each matched to its check in Appendix~\ref{app:lean}):
\begin{itemize}
\item $\ell$ is subadditive, invariant under conjugation and zero only at $1$;
\item $c^h=1$ and $\ell(c)=n>0$;
\item the exponent condition of Step~2;
\item the row shapes and the bound $b_{\max}$;
\item the row of the label of $c$ is $q^{n(m-1)}$, and every other row has a
factor $[m-1]_q$;
\item the identities at $m_1=1,\dots,n+1$ for $\nu\le n$, and those for
$\nu>n$ and $d\le b_{\max}$.
\end{itemize}

\emph{Conclusion:} for every $s\ge1$, every $k\in\Z/sh\Z$ of additive order
$d$, every label $X$ and every primitive $d$-th root of unity $\omega$ in a
field of characteristic $0$, the formula $\mathrm{rows}(X)$ at $m=sh+1$ is
regular at $\omega$, and its value there is the number of delta sequences
fixed by $\varphi^k$ whose bottom element has label $X$.

Instantiated with $G=W$, $\ell(w)=\operatorname{rank}(w-1)$ and
$\mathrm{cls}(w)$ the label of $W\cdot V^w$, this is
Theorem~\ref{thm:main}. Lean also proves that $\operatorname{rank}(w-1)$ on
a group of invertible matrices satisfies the three conditions on $\ell$.
The eleven theorems of the formal statement depend only on the axioms
\texttt{propext}, \texttt{Classical.choice} and \texttt{Quot.sound}.

Outside Lean remain: the finite computations themselves; that the checkers'
rational-function arithmetic computes the formal evaluation of a row, and
that their $\NC(W)$, order and length are the objects defined above, and
that their fixed-point polynomial is the $L_X$ of Step~2 (with $U_X$, the
chain counts and the block exponents as defined there), which a reader
confirms from the code; the identification of reflection length with
$\operatorname{rank}(w-1)$ \cite[Lemma~2]{Car}; the formulas of
\cite[\S3.6]{RS} as stated there, and that each row the checkers read is
the formula of the table; and the labeling of orbits of flats.
Appendix~\ref{app:lean} matches each hypothesis to its check.

\begin{appendices}
\section{The cases checked}\label{app:cases}

For each type the table gives the rank $n$, the Coxeter number $h$, the size
of $\NC(W)$, the number of orbits of flats, the number of pairs
$(h_1,\nu)$ with $h_1\mid h$, $1\le \nu\le n$ and $\gcd(h_1,\nu)=1$, the
number of pairs with $\nu>n$ and $\nu h/h_1\le b_{\max}$, $b_{\max}$, and the
largest $s$ for which checker B recounts fixed points by brute force, with
the literal action \eqref{eq:phi} and with $\rho$.

\begin{center}
\begin{tabular}{lrrrrrrrrr}
\toprule
Type & $n$ & $h$ & $|\NC(W)|$ & orbits & pairs, $\nu\le n$ & pairs, $\nu>n$ & $b_{\max}$ & $s$, \eqref{eq:phi} & $s$, $\rho$\\
\midrule
$G_2$ & 2 & 6 & 8 & 4 & 6 & 1 & 6 & 4 & 4\\
$F_4$ & 4 & 12 & 105 & 12 & 13 & 4 & 12 & 3 & 4\\
$E_6$ & 6 & 12 & 833 & 17 & 20 & 2 & 12 & 2 & 3\\
$E_7$ & 7 & 18 & 4160 & 32 & 27 & 4 & 18 & 1 & 2\\
$E_8$ & 8 & 30 & 25080 & 41 & 38 & 10 & 30 & 1 & 2\\
\bottomrule
\end{tabular}
\end{center}

\section{The formal statement and the checks}\label{app:lean}

The hypotheses of \texttt{RS16.Main.conjecture\_1\_4} and what establishes
each:
\begin{itemize}
\item subadditivity, conjugation invariance and $\ell(w)=0\iff w=1$ for
$\ell(w)=\operatorname{rank}(w-1)$: proved in Lean
(\texttt{reflection\_length});
\item $c^h=1$ and $\ell(c)=n$: a separate program computes the order of $c$
and the rank of $c-1$ in each type;
\item the exponent condition $h_1\mid \nu e_{j+1}+j$ for $0\le j<\nu$ (the
form of Step~2) for the block exponents:
proved in Lean for the recursion that checker A uses
(\texttt{checker\_exponents}), and such exponents exist for every coprime
pair (\texttt{exists\_exponents});
\item at most $n$ numerator and $n$ denominator brackets per row, the row
$q^{n(m-1)}$, the factor $[m-1]_q$ and the bound $b_{\max}$: checked on
every row by both checkers;
\item the identities for $\nu\le n$ and for $\nu>n$ with $d\le b_{\max}$:
checked by both checkers; a constant residue modulo the cyclotomic
polynomial gives the value at every primitive $d$-th root
(\texttt{checker\_value}).
\end{itemize}
Lean proves Step~2 (\texttt{step2\_fixed\_points},
\texttt{step2\_large\_m2}), Step~3 (\texttt{step3\_rows}), the vanishing of
rows with a factor $[m-1]_q$ (\texttt{rows\_with\_m1\_factor}) and Step~1
inside \texttt{conjecture\_1\_4}. Step~3 is proved by the finite differences
above, so it uses neither the positivity of $f_{e,1}$
\cite[Theorem~1.6(iv)]{RS} nor the check at $m_1=n+2$. In the Lean statement
a multichain is represented by its delta sequence
$(\delta_0;\delta_1,\dots,\delta_s)$, with $\delta_0=w_1$ as in
Section~\ref{sec:intro-main}.
\end{appendices}

\begin{thebibliography}{99}

\bibitem{AR}
C.~A.~Athanasiadis and V.~Reiner, \emph{Noncrossing partitions for the
group $D_n$}, SIAM J. Discrete Math. \textbf{18} (2004), 397--417;
\href{https://doi.org/10.1137/S0895480103432192}{doi:10.1137/S0895480103432192}.

\bibitem{Arm}
D.~Armstrong, \emph{Generalized noncrossing partitions and combinatorics of
Coxeter groups}, Mem. Amer. Math. Soc. \textbf{202} (2009), no.~949.

\bibitem{Bes}
D.~Bessis, \emph{The dual braid monoid}, Ann. Sci. École Norm. Sup. (4)
\textbf{36} (2003), 647--683;
\href{https://arxiv.org/abs/math/0101158}{arXiv:math/0101158}.

\bibitem{BR}
D.~Bessis and V.~Reiner, \emph{Cyclic sieving of noncrossing partitions for
complex reflection groups}, Ann. Comb. \textbf{15} (2011), 197--222;
\href{https://arxiv.org/abs/math/0701792}{arXiv:math/0701792}.

\bibitem{Car}
R.~W.~Carter, \emph{Conjugacy classes in the Weyl group}, Compositio Math.
\textbf{25} (1972), 1--59.

\bibitem{DJV}
T.~Douvropoulos and M.~Josuat-Vergès, \emph{The generalized cluster complex:
refined enumeration of faces and related parking spaces}, SIGMA \textbf{19}
(2023), 069; \href{https://doi.org/10.3842/SIGMA.2023.069}{doi:10.3842/SIGMA.2023.069}.

\bibitem{Kim}
D.~Kim, \emph{$q$-Kreweras numbers for coincidental Coxeter groups attached
to limit symbols}, \href{https://arxiv.org/abs/2012.08076}{arXiv:2012.08076}.

\bibitem{KM}
C.~Krattenthaler and T.~W.~Müller, \emph{Cyclic sieving for generalised
non-crossing partitions associated with complex reflection groups of
exceptional type}, in: I.~S.~Kotsireas and E.~V.~Zima (eds.), Advances in
Combinatorics: Waterloo Workshop in Computer Algebra, W80, Springer, Berlin,
2013, 209--247;
\href{https://doi.org/10.1007/978-3-642-30979-3_12}{doi:10.1007/978-3-642-30979-3\_12};
\href{https://arxiv.org/abs/1001.0028}{arXiv:1001.0028}; computational
details in \href{https://arxiv.org/abs/1001.0030}{arXiv:1001.0030}.

\bibitem{KS}
C.~Krattenthaler and C.~Stump, \emph{Positive $m$-divisible non-crossing
partitions and their Kreweras maps},
\href{https://arxiv.org/abs/2506.14996v1}{arXiv:2506.14996v1}.

\bibitem{Lean}
L.~de~Moura and S.~Ullrich, \emph{The Lean 4 theorem prover and programming
language}, in: Automated Deduction -- CADE 28, Lecture Notes in Comput. Sci.
\textbf{12699}, Springer, 2021, 625--635.

\bibitem{Mathlib}
The mathlib Community, \emph{The Lean mathematical library}, in: Proceedings
of the 9th ACM SIGPLAN International Conference on Certified Programs and
Proofs (CPP 2020), 367--381.

\bibitem{Pou}
L.~Pouillart, \emph{Cyclic sieving phenomena on parabolic classes of faces of
the cluster complex},
\href{https://arxiv.org/abs/2603.28242v1}{arXiv:2603.28242v1}.

\bibitem{RS}
V.~Reiner and E.~Sommers, \emph{Weyl group $q$-Kreweras numbers and cyclic
sieving}, Ann. Comb. \textbf{22} (2018), 819--874;
\href{https://arxiv.org/abs/1605.09172v2}{arXiv:1605.09172v2}.

\bibitem{Rho}
B.~Rhoades, \emph{Parking structures: Fuss analogs}, J. Algebraic Combin.
\textbf{40} (2014), 417--473;
\href{https://arxiv.org/abs/1205.4293}{arXiv:1205.4293}.

\bibitem{RSW}
V.~Reiner, D.~Stanton and D.~White, \emph{The cyclic sieving phenomenon},
J. Combin. Theory Ser. A \textbf{108} (2004), 17--50;
\href{https://doi.org/10.1016/j.jcta.2004.04.009}{doi:10.1016/j.jcta.2004.04.009}.

\end{thebibliography}
\end{document}
